% Section 6: Discussion (Course Specifics)
% What was easy, what was difficult, environmental impact

\section{Discussion}
\label{sec:discussion}

\subsection{What was Easy}
\label{sec:easy}

\textit{[To be completed based on implementation experience. Preliminary observations:]}

\begin{itemize}
    \item The bilinear layer architecture is straightforward to implement---it simply replaces nonlinearities with element-wise multiplication.
    \item The original paper's codebase was available and well-documented, providing clear reference for the eigendecomposition procedure.
    \item The MNIST dataset is small enough for rapid iteration during development.
\end{itemize}

\subsection{What was Difficult}
\label{sec:difficult}

\textit{[To be completed based on implementation experience. Preliminary observations:]}

\begin{itemize}
    \item Matching exact numerical values from the paper requires careful attention to random seed management across data loading, model initialization, and training.
    \item The eigendecomposition step is computationally expensive for large hidden dimensions.
    \item Determining the ``correct'' level of regularization that balances accuracy and interpretability requires hyperparameter search.
\end{itemize}

\subsection{Environmental Impact}
\label{sec:environmental}

Following responsible AI practices, we track and report the computational costs and CO2 emissions of our experiments using \texttt{codecarbon}.

% TODO: Uncomment after experiments
% \begin{table}[h]
%     \caption{Environmental impact of experiments}
%     \label{tab:environmental}
%     \centering
%     \begin{tabular}{lccc}
%         \toprule
%         Experiment & GPU Hours & Energy (kWh) & CO2 (kg) \\
%         \midrule
%         Phase 1: Reproduction (20 runs) & X.X & X.XX & X.XXX \\
%         Phase 2: Extensions (XX runs) & X.X & X.XX & X.XXX \\
%         Development \& debugging & X.X & X.XX & X.XXX \\
%         \midrule
%         \textbf{Total} & X.X & X.XX & X.XXX \\
%         \bottomrule
%     \end{tabular}
% \end{table}

\textit{[Table to be added: Environmental impact metrics from codecarbon]}

\textbf{Efficiency considerations:}
The CP-decomposition extension has potential environmental benefits.
By constraining the model to low rank structurally (rather than through regularization), we reduce the number of parameters and consequently:
\begin{itemize}
    \item Memory footprint during training and inference
    \item Computational cost of forward/backward passes
    \item Energy consumption at deployment scale
\end{itemize}

\subsection{Connection to FACT-AI Topics}
\label{sec:fact}

This work directly addresses \textbf{Transparency} (FACT-AI Topic 1.4):

\textbf{Weight-based vs. Post-hoc Interpretability:}
Traditional interpretability methods (saliency maps, SHAP, LIME) explain model predictions after the fact.
Bilinear MLPs offer a fundamentally different approach: the model's weights are \emph{directly} interpretable through eigendecomposition.
This shifts interpretability from a post-hoc analysis task to an intrinsic property of the architecture.

\textbf{Trade-offs:}
Our ablation study reveals the tension between accuracy and interpretability.
Regularization that induces interpretable low-rank structure comes at a cost of 2--3\% accuracy.
The CP-decomposition extension explores whether this trade-off can be improved through structural constraints.

\textbf{Implications for AI Accountability:}
If interpretable eigenvectors genuinely capture the model's decision-making features, they provide a basis for:
\begin{itemize}
    \item Auditing model behavior for biases
    \item Explaining individual predictions
    \item Identifying potential failure modes
\end{itemize}

\subsection{Communication with Original Authors}
\label{sec:authors}

\textit{[Document any communication with original authors, or note that no communication occurred.]}

\section{Conclusion}
\label{sec:conclusion}

\textit{[To be completed after experiments]}

We reproduced both main experimental sections from \citet{pearce2024bilinear}:
Section 4 (Vision), verifying that regularization induces low-rank interpretable structure in bilinear MLPs on MNIST;
and Section 5 (Language), analyzing bilinear attention in GPT-2 and comparing with Sparse Autoencoders.
Our ablation study [confirms/reveals] the relative importance of noise augmentation versus weight decay.
The CP-decomposition extension demonstrates that structural low-rank constraints [can/cannot] achieve comparable interpretability to regularization-induced low-rank.

\textbf{Limitations:}
\begin{itemize}
    \item The interpretability of eigenvectors is assessed qualitatively; future work could develop quantitative metrics.
    \item The language experiments focus on GPT-2; scaling to larger models remains future work.
\end{itemize}

\textbf{Future Work:}
\begin{itemize}
    \item Apply bilinear decomposition to larger-scale vision models
    \item Develop automated methods for assessing eigenvector interpretability
    \item Explore other tensor decompositions (Tucker, TT) as alternatives to CP
\end{itemize}
